\documentclass[11pt]{article}
\usepackage{amsmath,amssymb,amsthm}
\usepackage[margin=2.5cm]{geometry}

\newtheorem{theorem}{Theorem}
\newtheorem{lemma}[theorem]{Lemma}
\newtheorem{claim}[theorem]{Claim}
\theoremstyle{remark}
\newtheorem{remark}[theorem]{Remark}
\theoremstyle{definition}
\newtheorem*{statement}{Official statement}

\newcommand{\ip}[2]{\langle #1,#2\rangle}

\title{Problem 1, Cell 2: the orthogonality (chain) lemma}
\author{}
\date{}

\begin{document}
\maketitle

\begin{statement}
A line here always means a line through the origin of $\mathbb{R}^d$. The angle between two lines
$\ell,\ell'$ is the acute (non-obtuse) angle $\theta(\ell,\ell')\in[0,\pi/2]$ between them. If
$\ell,\ell'$ are spanned by unit vectors $x,x'$, then $\theta(\ell,\ell')=\arccos|\langle x,x'\rangle|$.
For lines $\ell_1,\dots,\ell_N$ in $\mathbb{R}^d$ (repetitions allowed), write
$S(\ell_1,\dots,\ell_N)=\sum_{1\le i<j\le N}\theta(\ell_i,\ell_j)$.

A statement about a chain of vectors, in which each vector may fail to be orthogonal only to its
immediate neighbours in the list. Let $m\ge2$ and let $x_1,\dots,x_m$ be unit vectors in
$\mathbb{R}^{m-1}$ with $\langle x_i,x_j\rangle=0$ whenever $|i-j|\ge2$. Prove that
$\sum_{i=1}^{m-1}\theta(x_i,x_{i+1})\le(m-2)\cdot\pi/2$, where $\theta(x,y)=\arccos|\langle x,y\rangle|$.
\end{statement}

\noindent Theorem~\ref{thm:main} below is exactly this statement: for unit vectors,
$\angle(x,y)=\theta(x,y)$.

\paragraph{Role in Problem~1.} This lemma is the chain case of the bound for $N=d+1$ lines in
$\mathbb{R}^d$ (the next cell of Problem~1). There, $d+1$ unit vectors in $\mathbb{R}^d$ are always
linearly dependent, and the lemma treats the dependent configurations whose non-orthogonal pairs
form a path. The cell asks for exactly this lemma.

\paragraph{Source.} The proof below is our own; we found no source for this lemma.

For nonzero $x,y\in\mathbb{R}^n$, the angle between the lines $\mathbb{R}x$ and $\mathbb{R}y$ is
\[
  \angle(x,y)=\arccos\frac{|\ip xy|}{\|x\|\,\|y\|}\in[0,\pi/2].
\]

\begin{theorem}[Orthogonality lemma]\label{thm:main}
Let $m\ge2$ and let $x_1,\dots,x_m\in\mathbb{R}^{m-1}$ be unit vectors with
\[
  \ip{x_i}{x_j}=0\qquad\text{whenever } |i-j|\ge2 .
\]
Put $\theta_i=\angle(x_i,x_{i+1})$ for $1\le i\le m-1$. Then
\[
  \sum_{i=1}^{m-1}\theta_i\le\frac{(m-2)\pi}{2}.
\]
\end{theorem}

\begin{remark}
Only $\dim\operatorname{span}(x_1,\dots,x_m)\le m-1$ is used. So the theorem also holds for unit
vectors in any $\mathbb{R}^n$ that span a subspace of dimension at most $m-1$.
\end{remark}

\section{Reductions}

\paragraph{Sign normalisation.}
Replacing some $x_i$ by $-x_i$ changes neither the lines $\mathbb{R}x_i$ (so neither the $\theta_i$)
nor the orthogonality hypotheses nor the ambient space. We choose signs one index at a time.
Keep $x_1$. For $i=1,2,\dots,m-1$ in turn, replace $x_{i+1}$ by $-x_{i+1}$ if
$\ip{x_i}{x_{i+1}}<0$. Step $i$ changes only $x_{i+1}$. So it does not affect
$\ip{x_{i'}}{x_{i'+1}}$ for $i'<i$, whose vectors have indices at most $i$ and are already fixed.
After step $m-1$,
\[
  c_i:=\ip{x_i}{x_{i+1}}\ge0\qquad(1\le i\le m-1).
\]
Since $\|x_i\|=\|x_{i+1}\|=1$, we get $c_i\in[0,1]$ and $\theta_i=\arccos c_i$.

\paragraph{Complementary angles.}
Put $\alpha_i=\frac\pi2-\theta_i\in[0,\frac\pi2]$. Then $c_i=\cos\theta_i=\sin\alpha_i$. Since
\[
  \sum_{i=1}^{m-1}\theta_i=\frac{(m-1)\pi}{2}-\sum_{i=1}^{m-1}\alpha_i,
\]
the theorem is equivalent to
\begin{equation}\label{eq:goal}
  \sum_{i=1}^{m-1}\alpha_i\ \ge\ \frac\pi2 .
\end{equation}

\section{Gram--Schmidt residuals}

Let $V_0=\{0\}$ and $V_j=\operatorname{span}(x_1,\dots,x_j)$ for $1\le j\le m$. Let $P_j$ be the orthogonal
projection onto $V_j$ (so $P_0=0$). Define
\[
  u_j=x_j-P_{j-1}x_j,\qquad r_j=\|u_j\|^2=\operatorname{dist}(x_j,V_{j-1})^2\qquad(1\le j\le m).
\]
Thus $u_j\perp V_{j-1}$ and, by Pythagoras, $r_j=1-\|P_{j-1}x_j\|^2$. Also $r_1=\|x_1\|^2=1$.
Finally, $r_j>0$ if and only if $x_j\notin V_{j-1}$.

For $1\le j\le m$ set $\beta_j=\alpha_1+\dots+\alpha_{j-1}$ (so $\beta_1=0$). Since every $\alpha_i\ge0$,
the sequence $\beta_j$ is nondecreasing.

\begin{lemma}[Recursion]\label{lem:rec}
Let $1\le j\le m-1$ and suppose $r_j>0$. Then $r_{j+1}=1-\dfrac{c_j^2}{r_j}$.
\end{lemma}

\begin{proof}
Since $u_j\in V_j$, $u_j\perp V_{j-1}$, and $x_j=P_{j-1}x_j+u_j$, we have the orthogonal
decomposition $V_j=V_{j-1}\oplus\mathbb{R}u_j$. Here $u_j\ne0$ because $r_j>0$. Hence
\[
  P_jy=P_{j-1}y+\frac{\ip{y}{u_j}}{r_j}\,u_j\qquad(y\in\mathbb{R}^{m-1}).
\]
Apply this to $y=x_{j+1}$. By hypothesis $x_{j+1}\perp x_1,\dots,x_{j-1}$ (their indices differ from
$j+1$ by at least $2$), so $x_{j+1}\perp V_{j-1}$. (For $j=1$ this is trivial, as $V_0=\{0\}$.)
Therefore $P_{j-1}x_{j+1}=0$. Moreover, since $P_{j-1}x_j\in V_{j-1}$,
\[
  \ip{x_{j+1}}{u_j}=\ip{x_{j+1}}{x_j}-\ip{x_{j+1}}{P_{j-1}x_j}=c_j-0=c_j .
\]
Hence $P_jx_{j+1}=\frac{c_j}{r_j}u_j$ and $\|P_jx_{j+1}\|^2=\frac{c_j^2}{r_j^2}\,r_j=\frac{c_j^2}{r_j}$.
Then $r_{j+1}=1-\|P_jx_{j+1}\|^2=1-c_j^2/r_j$.
\end{proof}

\begin{lemma}[Trigonometric inequality]\label{lem:trig}
Let $a,b\in[0,\pi/2]$ with $a+b<\pi/2$. Then
\[
  1-\frac{\sin^2a}{\cos^2b}\ \ge\ \cos^2(a+b).
\]
\end{lemma}

\begin{proof}
Note $\cos b>0$ because $b\le a+b<\pi/2$. The identity $\cos(x+y)\cos(x-y)=\cos^2x-\sin^2y$ gives
\[
  \cos(b+a)\cos(b-a)=\cos^2b\cos^2a-\sin^2b\sin^2a
  =\cos^2b(1-\sin^2a)-(1-\cos^2b)\sin^2a=\cos^2b-\sin^2a .
\]
Also $\cos(b-a)-\cos(b+a)=2\sin a\sin b\ge0$ because $a,b\in[0,\pi/2]$. Now $\cos(a+b)>0$, hence also $\cos(b-a)\ge\cos(a+b)>0$, and
$0<\cos^2b\le1$. Hence
\[
  1-\frac{\sin^2a}{\cos^2b}=\frac{\cos(a+b)\cos(b-a)}{\cos^2b}
  \ \ge\ \cos(a+b)\cos(b-a)\ \ge\ \cos(a+b)\cos(a+b).\qedhere
\]
\end{proof}

\begin{claim}\label{cl:ind}
For every $j\in\{1,\dots,m\}$ with $\beta_j<\pi/2$ we have $r_j\ge\cos^2\beta_j>0$.
\end{claim}

\begin{proof}
Induction on $j$. For $j=1$: $\beta_1=0$ and $r_1=1=\cos^20$.

Step $j\to j+1$ ($1\le j\le m-1$). Assume $\beta_{j+1}=\beta_j+\alpha_j<\pi/2$. Then
$\beta_j\le\beta_{j+1}<\pi/2$, so by the induction hypothesis $r_j\ge\cos^2\beta_j>0$. By
Lemma~\ref{lem:rec}, $r_{j+1}=1-c_j^2/r_j$. Using $c_j^2=\sin^2\alpha_j$ and
$r_j\ge\cos^2\beta_j>0$, we get $c_j^2/r_j\le\sin^2\alpha_j/\cos^2\beta_j$. Hence, by
Lemma~\ref{lem:trig} with $a=\alpha_j$ and $b=\beta_j$ (both lie in $[0,\pi/2]$, and
$a+b=\beta_{j+1}<\pi/2$),
\[
  r_{j+1}\ \ge\ 1-\frac{\sin^2\alpha_j}{\cos^2\beta_j}\ \ge\ \cos^2(\alpha_j+\beta_j)=\cos^2\beta_{j+1}>0,
\]
where the last inequality holds because $0\le\beta_{j+1}<\pi/2$.
\end{proof}

\section{Proof of Theorem~\ref{thm:main}}

Suppose, for contradiction, that \eqref{eq:goal} fails, i.e.\ $\beta_m=\sum_{i=1}^{m-1}\alpha_i<\pi/2$.
Because $(\beta_j)$ is nondecreasing, $\beta_j<\pi/2$ for every $j=1,\dots,m$. By Claim~\ref{cl:ind},
$r_j>0$ for all $j$, i.e.\ $x_j\notin\operatorname{span}(x_1,\dots,x_{j-1})$ for all $j$. Then
$x_1,\dots,x_m$ are linearly independent. (If $\sum\lambda_ix_i=0$ with not all $\lambda_i=0$, let
$j$ be the largest index with $\lambda_j\ne0$; then $x_j\in V_{j-1}$, a contradiction.) This gives
$m$ linearly independent vectors in $\mathbb{R}^{m-1}$, which is impossible. Hence
$\sum\alpha_i\ge\pi/2$, which is \eqref{eq:goal}, and so $\sum\theta_i\le(m-2)\pi/2$. \qed

\section{Edge cases}

\begin{itemize}
\item $m=2$: the vectors lie in $\mathbb{R}^1$, so $x_2=\pm x_1$ and $\theta_1=0=(m-2)\pi/2$. The proof
  covers this case: $\alpha_1<\pi/2$ would give $r_2=1-\sin^2\alpha_1=\cos^2\alpha_1>0$, i.e.\ two
  independent vectors in $\mathbb{R}^1$. There are no orthogonality hypotheses when $m=2$.
\item $c_i=1$ for some $i$ (i.e.\ $x_{i+1}=x_i$ after normalisation): then $\alpha_i=\pi/2$ and
  \eqref{eq:goal} holds trivially. In the argument, $\beta_j\ge\pi/2$ for $j>i$, so the claim is
  never invoked beyond index $i$. In particular we never divide by an $r_j$ that might vanish.
\item $c_i=0$: then $\alpha_i=0$, and the recursion gives $r_{i+1}=1$, consistent with
  $x_{i+1}\perp V_i$.
\item Division by $r_j$ occurs only in Lemma~\ref{lem:rec}, and only under the hypothesis $r_j>0$.
  In Claim~\ref{cl:ind} that hypothesis comes from the induction.
\end{itemize}

\begin{remark}[Sharpness]
Equality is attained. For example, take $x_1=e_1$, $x_2=e_2,\dots,x_{m-1}=e_{m-1}$ and
$x_m=e_{m-1}$ in $\mathbb{R}^{m-1}$. Then $\theta_1=\dots=\theta_{m-2}=\pi/2$ and $\theta_{m-1}=0$.
For $m=3$, the chain $e_1$, $(\cos t,\sin t)$, $e_2$ in $\mathbb{R}^2$ with $t\in[0,\pi/2]$ gives
equality for every $t$.
\end{remark}

\end{document}
